% ============================================================
% 5_E_admin_teacher_ui.tex — 管理员端 + 教师端 UI 详细设计
% 编写人：E  日期：2026-07-15
% 职责：管理员页面控件 / 教师批阅交互 / 前端校验规则 / 状态管理
% ============================================================

\section{管理员端页面详细设计}

\subsection{页面1：仪表盘页（AdminDashboard.vue）}

\subsubsection{页面布局}
仪表盘页采用四卡片+双图表的网格布局，上方为四个指标卡片排列成一行，下方为两个图表并排展示。页面顶部右侧提供时间范围选择器（el-date-picker），筛选后全局刷新图表数据。

\subsubsection{控件清单}
\begin{table}[htbp]
\centering
\caption{仪表盘页控件明细}
\label{tab:admin_dashboard}
\begin{tabular}{|c|c|c|c|}
\hline
\textbf{控件类型} & \textbf{用途} & \textbf{数据来源} & \textbf{交互逻辑} \\
\hline
StatCard（自定义） & 显示统计数字 & GET /api/admin/stats & 页面加载时一次性获取 \\
\hline
el-date-picker & 时间范围筛选 & — & 选择触发图表数据刷新 \\
\hline
Chart（ECharts） & 论文状态分布饼图 & GET /api/admin/stats/papers-by-status & 按选中时间范围查询 \\
\hline
Chart（ECharts） & 导出趋势折线图 & GET /api/admin/stats/export-trend & 按选中时间范围查询 \\
\hline
\end{tabular}
\end{table}

\subsubsection{数据接口}
\begin{lstlisting}[language=Java,caption={仪表盘数据接口定义}]
// 获取统计数据概览
GET /api/admin/stats
Response: {
  totalUsers: 150,      // 总用户数
  totalStudents: 120,   // 学生数
  totalTeachers: 25,    // 教师数
  totalAdmins: 5,       // 管理员数
  totalPapers: 200,     // 论文总数
  totalTemplates: 10,   // 模板总数
  todayExports: 15      // 今日导出数
}

// 获取论文状态分布
GET /api/admin/stats/papers-by-status?start=&end=
Response: [{ status: "DRAFT", count: 80 }, ...]

// 获取导出趋势
GET /api/admin/stats/export-trend?start=&end=
Response: [{ date: "2026-07-10", count: 5 }, ...]
\end{lstlisting}

\subsection{页面2：学院管理页（CollegeAdmin.vue）}

\subsubsection{控件清单}
\begin{table}[htbp]
\centering
\caption{学院管理页控件明细}
\begin{tabular}{|c|c|c|c|}
\hline
\textbf{控件类型} & \textbf{字段/用途} & \textbf{校验规则} & \textbf{数据操作} \\
\hline
el-table & 学院列表展示 & — & 行数据由API填充 \\
\hline
el-input & 搜索框 & 最大长度50 & 输入触发前端过滤 \\
\hline
el-dialog + el-form & 新增/编辑表单 & 名称非空、代码唯一 & POST/PUT API \\
\hline
el-button & 删除操作 & — & 二次确认后 DELETE \\
\hline
\end{tabular}
\end{table}

\subsubsection{交互逻辑}
\begin{enumerate}
    \item 页面加载时调用 GET /api/admin/colleges 获取学院列表，渲染至el-table；
    \item 点击"新增学院"按钮，弹出Dialog，表单包含name（必填）、code（必填，唯一）两字段；
    \item 表单校验通过后提交 POST /api/admin/colleges，成功后关闭Dialog并刷新列表；
    \item 点击编辑图标，弹出Dialog预填当前行数据，提交 PUT /api/admin/colleges/{id}；
    \item 点击删除图标，弹出二次确认框，确认后调用 DELETE /api/admin/colleges/{id}，若服务端返回关联校验错误（如存在关联模板），在前端以 ElMessage.error 展示错误信息。
\end{enumerate}

\subsection{页面3：模板管理页（TemplateAdmin.vue）}

\subsubsection{页面布局}
双栏布局：左侧为模板卡片列表（宽度320px），右侧为模板编辑面板（自适应宽度）。左侧上方提供筛选栏（按类型下拉选择 + 按学院下拉选择 + 搜索框）。右侧编辑面板使用Tabs组织三个配置子页面。

\subsubsection{模板卡片组件设计}
每个模板卡片展示：模板名称（加粗16px）、类型标签（el-tag，颜色按类型区分：毕业=蓝色、课程=绿色、项目=橙色）、当前版本号、启用状态开关（el-switch，独立提交 PUT /api/admin/templates/{id}/toggle）。卡片选中状态以左侧蓝色边框高亮标记。

\subsubsection{封面字段配置面板}
封面字段配置采用拖拽排序的列表实现：
\begin{itemize}
    \item 列表项显示字段标签（如"论文题目"）、是否必填标记、字段类型（文本/日期/下拉选择）；
    \item 每个列表项右侧提供编辑和删除按钮；
    \item 列表底部有"添加字段"按钮，点击弹出字段配置Dialog（字段标签、是否必填、字段类型、默认值）；
    \item 拖拽使用vuedraggable库实现，拖拽完成后自动更新cover_fields数组顺序，标记为未保存状态。
\end{itemize}

\subsubsection{章节结构配置面板}
使用el-tree组件展示树形章节结构，每个节点右侧显示章节类型标识（chapter/section/subsection）。根节点（论文标题）不可删除。交互功能包括：
\begin{itemize}
    \item 右键菜单（@node-contextmenu事件）：新增同级章节、新增子章节、重命名、删除、上移、下移；
    \item 拖拽（draggable属性）：跨层级拖拽调整章节归属关系，拖拽时校验不产生循环引用；
    \item 节点点击选中后，右侧展示该章节的详细配置（章节标题输入框、章节类型下拉选择、默认内容占位符）。
\end{itemize}

\subsubsection{格式参数配置面板}
格式参数配置采用el-form垂直布局，分组排列：
\begin{itemize}
    \item 页面设置组：纸张大小下拉选择（A4/Legal/Letter）、上/下/左/右边距数字输入框（el-input-number，单位cm，步进0.1）；
    \item 字体设置组：正文字体下拉选择、正文字号下拉选择、一级标题字体/字号、二级标题字体/字号、三级标题字体/字号；
    \item 段落设置组：行距倍数滑块（el-slider，范围1.0-2.5，步进0.1）、段前间距数字输入框、段后间距数字输入框；
    \item 页眉页脚组：奇偶页眉内容输入框、页码格式下拉选择、页码位置单选组。
\end{itemize}

配置变更时，顶部的"保存"按钮由禁用状态变为可用，底部显示"有未保存的更改"提示条。保存时调用 PUT /api/admin/templates/{id}/config，后端自动执行版本递增逻辑。

\section{教师端页面详细设计}

\subsection{页面1：批阅任务列表页（ReviewList.vue）}

\subsubsection{页面布局}
顶部为筛选操作栏，主体为表格视图。筛选栏包含：学院下拉选择框、论文类型下拉选择框、状态选项卡（全部/待批阅/已批阅/已退回）。

\subsubsection{控件清单}
\begin{table}[htbp]
\centering
\caption{批阅任务列表页控件明细}
\begin{tabular}{|c|c|c|c|}
\hline
\textbf{控件类型} & \textbf{用途} & \textbf{数据来源} & \textbf{交互} \\
\hline
el-tabs & 状态筛选 & — & 切换触发数据重新加载 \\
\hline
el-select & 学院筛选 & GET /api/colleges & 联动刷新列表 \\
\hline
el-table & 论文列表展示 & GET /api/teacher/reviews & 分页加载 \\
\hline
el-tag & 状态标签 & 状态字段映射 & 颜色区分：橙=待批阅、绿=已批阅、灰=已退回 \\
\hline
el-button & 进入批阅 & — & 点击跳转 /teacher/reviews/{id} \\
\hline
\end{tabular}
\end{table}

\subsubsection{数据接口}
\begin{lstlisting}[language=Java,caption={批阅列表接口定义}]
// 获取批阅任务列表
GET /api/teacher/reviews?status=&collegeId=&page=&size=
Response: {
  records: [{
    id: 1,
    thesisId: 10,
    thesisTitle: "基于Spring Boot的在线论文生成系统",
    studentName: "张三",
    collegeName: "计算机科学与技术学院",
    submittedAt: "2026-07-14 10:30:00",
    status: "SUBMITTED",   // SUBMITTED / REVIEWED / RETURNED
    annotationCount: 0,    // 批注数量
    score: null,           // 分数
    grade: null            // 等级
  }],
  total: 50,
  page: 1,
  size: 20
}
\end{lstlisting}

\subsection{页面2：论文批阅页（ReviewDetail.vue）}

\subsubsection{页面布局}
三栏布局。左侧栏宽度260px展示章节树导航，中间栏自适应宽度展示论文全文预览，右侧栏宽度340px展示批注管理面板。顶部操作栏包含论文标题、学生信息、操作按钮区（评阅完成按钮、返回列表按钮）。

\subsubsection{章节树面板}
使用el-tree组件展示论文完整章节树，章节节点显示标题和该章节的批注数量角标。当前选中的章节高亮显示。教师点击章节节点时切换中间栏的预览内容，同时右侧批注面板加载对应章节的批注列表。

\subsubsection{论文预览区}
论文内容以HTML渲染展示，容器样式模拟A4纸张效果（白色背景、阴影边框、标准页边距）。核心交互逻辑：
\begin{itemize}
    \item 使用contenteditable=false的div容器展示HTML内容，确保教师不可编辑；
    \item 已批注的文字以黄色背景色（\#FFF5CC）高亮标记，并添加波浪下划线样式；
    \item 教师选中文字时触发selectionchange事件，获取选中文本的起始偏移量和长度；
    \item 选中文字后在选中区域上方浮动出现"添加批注（Alt+A）"按钮，点击后触发右侧批注编辑区。
\end{itemize}

\subsubsection{批注定位算法}

教师选中文段时，需要精确记录选中文字在原始HTML内容中的位置，以便将来再次打开批阅时能正确定位到同一段文字。

\begin{lstlisting}[language=JavaScript,caption={批注位置定位算法}]
// 获取选中文字在HTML纯文本中的偏移量
function getSelectionOffset(container) {
    const selection = window.getSelection();
    if (!selection.rangeCount) return null;

    const range = selection.getRangeAt(0);
    const preRange = document.createRange();
    preRange.selectNodeContents(container);
    preRange.setEnd(range.startContainer, range.startOffset);

    // 计算纯文本起始偏移（忽略HTML标签）
    const startOffset = preRange.toString().length;
    const textLength = selection.toString().length;

    return { startOffset, textLength, selectedText: selection.toString() };
}

// 打开批阅时定位到已有批注位置
function scrollToAnnotation(annotation) {
    const container = document.getElementById('preview-content');
    const text = container.textContent || container.innerText;
    const beforeText = text.substring(0, annotation.startOffset);
    const targetText = text.substring(
        annotation.startOffset,
        annotation.startOffset + annotation.textLength
    );

    // 使用 Range API 定位并高亮
    const range = document.createRange();
    // 遍历文本节点查找目标位置（省略具体遍历实现）
    // 定位后滚动到可见区域
    range.collapse(true);
    const rect = range.getBoundingClientRect();
    container.scrollTop += rect.top - 200; // 偏移200px保证上下文可见
}
\end{lstlisting}

\subsubsection{批注面板}
右侧面板分为上下两部分：
\begin{itemize}
    \item \textbf{上部 — 批注列表}：按创建时间降序排列当前章节的全部批注。每条批注展示：批注内容摘要（截取前50字）、被选中原文（灰色斜体，用引号包裹）、创建时间（相对时间如"3分钟前"）。点击批注项时，中间栏自动滚动定位到对应原文位置；
    \item \textbf{下部 — 批注编辑区}：使用textarea组件展示教师输入的批注内容。当教师选中原文时，自动填充"选中文字"预览区域（只读），textarea获取焦点等待输入。提供"提交批注"按钮（POST /api/annotations）和"取消"按钮。
\end{itemize}

\subsection{评分对话框详细设计}

\subsubsection{控件布局}
评分对话框使用el-dialog组件，宽度620px，标题为"评阅论文"。内容表单分三部分：
\begin{enumerate}
    \item \textbf{评语区域}：使用Tiptap轻量级富文本编辑器（或el-input type="textarea"），placeholder为"请输入整体评语（可选）"，最大长度5000字符；
    \item \textbf{评分区域}：el-slider水平滑块，最小值0，最大值100，步进1，显示当前分值（大号数字：48px字体），评分下方自动展示对应等级标签（el-tag）：
    \begin{itemize}
        \item $\ge$90：红色标签"优秀"
        \item 80-89：蓝色标签"良好"
        \item 70-79：黄色标签"中等"
        \item 60-69：绿色标签"及格"
        \item <60：灰色标签"不及格"
    \end{itemize}
    \item \textbf{操作区域}：el-radio-group横向排列，"通过（REVIEWED）"为绿色按钮样式，"退回（RETURNED）"为红色按钮样式。选择"退回"后下方动态展开textarea输入框用于填写退回原因（必填，最小10字符）。
\end{enumerate}

\subsubsection{校验规则}
\begin{table}[htbp]
\centering
\caption{评分对话框校验规则}
\begin{tabular}{|c|c|c|c|}
\hline
\textbf{字段} & \textbf{规则} & \textbf{错误提示} & \textbf{触发时机} \\
\hline
评语 & 最大5000字符 & 超出长度限制 & 输入时 \\
\hline
评分 & 0-100整数 & 请输入0-100之间的整数 & 提交时 \\
\hline
操作类型 & 必选 & 请选择操作类型 & 提交时 \\
\hline
退回原因 & 操作类型=RETURNED时必填，最小10字符 & 退回原因至少10个字符 & 提交时 \\
\hline
\end{tabular}
\end{table}

\subsubsection{提交流程}
\begin{enumerate}
    \item 教师点击Dialog的"确认提交"按钮；
    \item 前端执行校验：评分范围、退回原因条件必填；
    \item 校验通过后调用 POST /api/teacher/reviews，请求体包含 thesisId、commentHtml、score、action、returnReason（可选）；
    \item 提交成功后关闭Dialog，前端显示 ElMessage.success("评阅完成")，跳转回批阅任务列表页；
    \item 提交失败时显示 ElMessage.error，保留当前Dialog中的表单数据供教师修正。
\end{enumerate}

\subsection{前端状态管理}

教师端批阅页面涉及多个组件间的状态共享，使用Pinia store进行管理：

\begin{lstlisting}[language=JavaScript,caption={批阅页面Pinia Store定义}]
// stores/review.ts
export const useReviewStore = defineStore('review', () => {
    // 当前批阅的论文信息
    const currentThesis = ref(null)
    // 章节树列表
    const sections = ref([])
    // 当前选中的章节ID
    const activeSectionId = ref(null)
    // 当前章节的全部批注
    const annotations = ref([])
    // 当前选中的文本信息（用于添加批注）
    const selectedText = ref(null)
    // 加载状态
    const loading = ref(false)

    // 加载论文批阅数据
    async function loadReview(thesisId) {
        loading.value = true
        const [thesisRes, sectionsRes] = await Promise.all([
            api.getThesisDetail(thesisId),
            api.getSections(thesisId)
        ])
        currentThesis.value = thesisRes.data
        sections.value = sectionsRes.data
        loading.value = false
    }

    // 切换章节时更新批注列表
    async function switchSection(sectionId) {
        activeSectionId.value = sectionId
        const res = await api.getAnnotations(currentThesis.value.id, sectionId)
        annotations.value = res.data
    }

    return {
        currentThesis, sections, activeSectionId,
        annotations, selectedText, loading,
        loadReview, switchSection
    }
})
\end{lstlisting}
